\chapter{Μεθοδολογία Αξιολόγησης}

\section{Εισαγωγή στη Μεθοδολογία}
Η αξιολόγηση της αναπτυχθείσας πλατφόρμας σχεδιάστηκε με στόχο να τεκμηριωθεί τόσο η τεχνική της αποτελεσματικότητα όσο και η πρακτική της αξία σε πραγματικό σχολικό περιβάλλον. Η μεθοδολογία εφαρμόστηκε στο πεδίο μελέτης των Εκπαιδευτηρίων Πλάτων και οργανώθηκε γύρω από τρεις διακριτούς, αλλά αλληλοσυμπληρούμενους, άξονες: (α) την αξιολόγηση του μαθησιακού περιβάλλοντος μέσω της ακουστικής άνεσης, (β) την αξιολόγηση της ενεργειακής απόδοσης με κανονικοποίηση των καταναλώσεων σε \textlatin{kWh}/ημέρα, και (γ) την τεχνοοικονομική αξιολόγηση της επένδυσης με χρήση δεικτών \textlatin{CAPEX}, \textlatin{OPEX savings} και πενταετούς \textlatin{ROI}.

Η μεθοδολογική προσέγγιση ακολούθησε δύο διαδοχικά στάδια υλοποίησης. Στο πρώτο στάδιο, το οποίο αντιστοιχεί στο τοπικό πρωτότυπο (\textlatin{Minimum Viable Product - MVP}), η συλλογή του \textlatin{baseline} και των πρώτων τηλεμετρικών δεδομένων πραγματοποιήθηκε μέσω της υποδομής \textlatin{openHAB}/\textlatin{InfluxDB} στο φυσικό πεδίο. Το στάδιο αυτό επέτρεψε την αρχική επικύρωση των αισθητήρων, των κανόνων ελέγχου και της ροής δεδομένων από το επίπεδο των συσκευών προς το επίπεδο εποπτείας. Στο δεύτερο στάδιο, η πλήρης αναλυτική επεξεργασία, οι προβολές εξοικονόμησης, οι τεχνοοικονομικοί δείκτες και τα εργαλεία \textlatin{AI Energy Forecasting} ενσωματώθηκαν στη νέα \textlatin{Enterprise Cloud SaaS} πλατφόρμα, η οποία παρείχε πιο ώριμο περιβάλλον ανάλυσης, ιστορικής σύγκρισης και λήψης αποφάσεων για τον τελικό χρήστη.

Η παραπάνω διάκριση είναι ουσιώδης για την ερμηνεία των αποτελεσμάτων. Το τοπικό \textlatin{MVP} λειτούργησε ως πλαίσιο επιχειρησιακής μέτρησης και αρχικής απόδειξης σκοπιμότητας, ενώ η μετέπειτα \textlatin{SaaS} υλοποίηση μετέφερε την αξιολόγηση σε ένα ανώτερο επίπεδο, όπου η συσχέτιση δεδομένων, η εξαγωγή δεικτών απόδοσης και η προβολή οικονομικών ωφελειών μπορούσαν να εκτελούνται με περισσότερο τυποποιημένο και επεκτάσιμο τρόπο.

Η διασφάλιση της ποιότητας των μετρήσεων αποτέλεσε αναπόσπαστο μέρος της μεθοδολογίας. Πριν την έναρξη της συλλογής δεδομένων, πραγματοποιήθηκε βαθμονόμηση όλων των αισθητήρων έναντι πιστοποιημένων οργάνων αναφοράς. Ειδικότερα, οι κόμβοι ακουστικής παρακολούθησης (\textlatin{ESP32/INMP441}) βαθμονομήθηκαν έναντι βιομηχανικού ηχομέτρου, επιτυγχάνοντας ακρίβεια ±2 \textlatin{dBA}. Οι αισθητήρες θερμοκρασίας και υγρασίας (\textlatin{DHT22}) βαθμονομήθηκαν με αποτέλεσμα ακρίβεια ±0.5°C και ±2\% \textlatin{RH} αντίστοιχα. Η συνεχής παρακολούθηση της ποιότητας δεδομένων μέσω αυτοματοποιημένων ρουτινών ανίχνευσης εκτός ορίων (\textlatin{out-of-range detection}) διασφάλισε \textlatin{data quality score} 98.5\% καθ' όλη την περίοδο αξιολόγησης.

\section{Αξιολόγηση Μαθησιακού Περιβάλλοντος (Ακουστική Άνεση)}
Η αξιολόγηση του μαθησιακού περιβάλλοντος βασίστηκε στην παρακολούθηση της ακουστικής συμπεριφοράς των αιθουσών διδασκαλίας, με έμφαση στη λειτουργική επίδραση του κλειστού βρόχου ελέγχου στο επίπεδο θορύβου. Ως επιχειρησιακό κατώφλι παρέμβασης ορίστηκε η τιμή των 55 \textlatin{dBA} (\textlatin{LAeq}) κατά τη διάρκεια ενεργού διδακτικής ώρας. Η τιμή αυτή δεν αποτελεί θεσμοθετημένο όριο για ενεργή τάξη, αλλά επιλέχθηκε ως επιχειρησιακός συμβιβασμός μεταξύ της τιμής αναφοράς \textlatin{WHO} για άδεια αίθουσα ($\leq$35 \textlatin{dBA}) \cite{WHO1999, Kephalopoulos2015} και της πρακτικής ακουστικής δυναμικής κατά τη διδασκαλία, σύμφωνα με τη μεθοδολογία παθητικής αυτορρύθμισης συμπεριφοράς (\textlatin{behavioral nudging}) που προτείνεται στη βιβλιογραφία \cite{MogasRecalde2021}.

Στην πρώτη φάση αξιολόγησης, το σύστημα λειτουργούσε ως παθητικός παρατηρητής, καταγράφοντας το πλήθος και τη συχνότητα των υπερβάσεων του ορίου χωρίς να παράγει φυσική παρέμβαση στον χώρο. Η περίοδος αυτή παρείχε το σημείο αναφοράς για την κατανόηση της βασικής ακουστικής συμπεριφοράς των αιθουσών. Στη δεύτερη φάση, ενεργοποιήθηκε ο αμφίδρομος έλεγχος, κατά τον οποίο οι συνεχείς υπερβάσεις του ορίου οδηγούσαν στην ενεργοποίηση της διακριτικής οπτικής ένδειξης (\textlatin{quiet-lamp}). Με τον τρόπο αυτό, το σύστημα δεν περιοριζόταν στην καταγραφή, αλλά παρήγαγε μια ήπια μορφή ανατροφοδότησης που μπορούσε να επηρεάσει θετικά τη συμπεριφορά των μαθητών και τη συνολική λειτουργία της τάξης.

Η αξιολόγηση δεν στηρίχθηκε μόνο στην παρατήρηση μεμονωμένων μετρήσεων, αλλά στο πλήθος των υπερβάσεων ανά διδακτική περίοδο, στη χρονική διάρκειά τους και στη συγκριτική αποτίμηση των δύο φάσεων λειτουργίας. Οι παράμετροι και τα όρια που χρησιμοποιήθηκαν για την ακουστική αξιολόγηση συνοψίζονται στον Πίνακα~\ref{tab:noise_limits}. Η αρχική συλλογή των δεδομένων πραγματοποιήθηκε στο τοπικό \textlatin{MVP}, ενώ η μεταγενέστερη ανάλυση τάσεων και η οπτικοποίηση των αποτελεσμάτων μεταφέρθηκαν στο περιβάλλον της \textlatin{Cloud SaaS} πλατφόρμας, όπου μπορούσαν να παραχθούν συγκεντρωτικοί δείκτες και περισσότερο φιλικές προς τον τελικό χρήστη αναφορές.

\begin{table}[htbp]
\centering
\caption{Όρια θορύβου και παράμετροι ενεργοποίησης οπτικής ένδειξης}
\label{tab:noise_limits}
\begin{tabular}{|p{4.8cm}|p{3cm}|p{6.2cm}|}
\hline
\textbf{Παράμετρος} & \textbf{Τιμή} & \textbf{Ρόλος / Περιγραφή} \\
\hline
Συνιστώμενο επίπεδο υποβάθρου (Άδεια τάξη) & $\leq$ 35 \textlatin{dBA} & Τιμή αναφοράς για ήσυχο σχολικό περιβάλλον χωρίς ενεργή διδασκαλία. \\
\hline
Όριο ενεργοποίησης ένδειξης (Κατά τη διδασκαλία) & 55 \textlatin{dBA} & Επιχειρησιακό κατώφλι παρέμβασης για την ενεργοποίηση της οπτικής ανατροφοδότησης κατά την ενεργή διδακτική ώρα. \\
\hline
Χρονικό παράθυρο επιβεβαίωσης (\textlatin{Dwell Time}) & 20 δευτερόλεπτα & Συνεχής υπέρβαση του ορίου για το ελάχιστο αυτό διάστημα απαιτείται πριν ενεργοποιηθεί η ένδειξη, ώστε να αποφεύγονται στιγμιαίες ψευδο-ενεργοποιήσεις. \\
\hline
Περίοδος εφαρμογής & Ενεργή διδακτική ώρα & Ο κανόνας ελέγχου εφαρμόζεται μόνο στη διάρκεια μαθήματος και όχι σε διαλείμματα ή εκτός ωρολογίου προγράμματος. \\
\hline
Ενεργοποιητής (\textlatin{Actuator}) & Οπτική λυχνία (\textlatin{quiet-lamp}) & Μη παρεμβατικός μηχανισμός ανατροφοδότησης που στοχεύει στην ήπια αυτορρύθμιση της συμπεριφοράς των μαθητών. \\
\hline
\end{tabular}
\end{table}

\section{Μεθοδολογία Ενεργειακής Αξιολόγησης}
Η αξιολόγηση της ενεργειακής επίδρασης της πλατφόρμας βασίστηκε στη σύγκριση της κατανάλωσης ηλεκτρικής ενέργειας πριν και μετά την επιχειρησιακή εφαρμογή του ψηφιακού διδύμου. Ως πρωτογενές υλικό χρησιμοποιήθηκαν οι επίσημοι λογαριασμοί ηλεκτρικής ενέργειας της εγκατάστασης για την περίοδο από τον Σεπτέμβριο του 2024 έως τον Νοέμβριο του 2025. Επειδή οι επιμέρους περίοδοι τιμολόγησης δεν είχαν σταθερή διάρκεια σε ημέρες, η σύγκριση δεν θα μπορούσε να γίνει αξιόπιστα σε επίπεδο απόλυτων μηνιαίων τιμών. Για τον λόγο αυτό εφαρμόστηκε κανονικοποίηση των καταναλώσεων σε \textlatin{kWh}/ημέρα, καθώς οι περίοδοι τιμολόγησης κυμαίνονταν μεταξύ 28 και 63 ημερών, καθιστώντας αδύνατη την άμεση σύγκριση απόλυτων τιμών, ώστε να προκύψει μια συγκρίσιμη και μεθοδολογικά ομοιογενής βάση αξιολόγησης.

Η ενεργειακή μεθοδολογία υλοποιήθηκε επίσης σε δύο στάδια. Στο πρώτο στάδιο, το τοπικό \textlatin{MVP} κατέγραψε το \textlatin{baseline} της εγκατάστασης και τα πρώτα επιχειρησιακά δεδομένα μέσω της σύζευξης \textlatin{openHAB}, \textlatin{MQTT} και \textlatin{InfluxDB}. Το στάδιο αυτό λειτούργησε ως σημείο αναφοράς για τη συμπεριφορά της εγκατάστασης πριν από την πλήρη ωρίμανση των αυτοματισμών. Στο δεύτερο στάδιο, τα ιστορικά δεδομένα και οι λογικές ελέγχου μεταφέρθηκαν στη νέα \textlatin{Enterprise Cloud SaaS} πλατφόρμα, όπου κατέστη δυνατή η ενοποιημένη ανάλυση χρονοσειρών, η σύγκριση \textlatin{baseline} και βελτιστοποιημένης λειτουργίας, καθώς και η παραγωγή προβλέψεων μέσω \textlatin{AI Energy Forecasting}.

Για την ερμηνεία των αποτελεσμάτων διαχωρίστηκαν δύο επίπεδα σύγκρισης. Το πρώτο αφορά τη γενική ακατέργαστη σύγκριση (\textlatin{raw comparison}) του συνόλου των διαθέσιμων περιόδων πριν και μετά. Το δεύτερο αφορά τη σύγκριση της περιόδου βάσης (\textlatin{baseline}) με την περίοδο πλήρους επιχειρησιακής βελτιστοποίησης, κατά την οποία οι κανόνες ελέγχου είχαν ενεργοποιηθεί σταθερά και η πλατφόρμα λειτουργούσε με ώριμο τρόπο. Η δεύτερη προσέγγιση χρησιμοποιήθηκε ως κύριο λειτουργικό \textlatin{KPI}, καθώς αντανακλά με μεγαλύτερη πιστότητα την απόδοση του συστήματος υπό συνθήκες πλήρους λειτουργίας.

Η νέα \textlatin{Cloud SaaS} πλατφόρμα επέτρεψε, επιπλέον, την ενσωμάτωση αναλυτικών εργαλείων πρόβλεψης κόστους και μελλοντικής κατανάλωσης. Η δυνατότητα αυτή δεν αναιρεί τη μεθοδολογία της αρχικής συλλογής, αλλά τη συμπληρώνει, παρέχοντας ένα πιο ώριμο περιβάλλον αξιολόγησης για διοικητικούς και τεχνικούς χρήστες, με καλύτερη αναγνωσιμότητα και μεγαλύτερη επιχειρησιακή αξία.

\section{\texorpdfstring{Τεχνοοικονομική Μεθοδολογία (Υπολογισμός \textlatin{ROI})}{Τεχνοοικονομική Μεθοδολογία (Υπολογισμός ROI)}}
Η τεχνοοικονομική αξιολόγηση σχεδιάστηκε με στόχο να εκτιμηθεί κατά πόσο η αναπτυχθείσα πλατφόρμα μπορεί να υποστηρίξει μια βιώσιμη επένδυση για ένα σχολικό συγκρότημα. Η ανάλυση στηρίχθηκε στη σύγκριση του αρχικού κόστους κεφαλαίου (\textlatin{CAPEX}) με τα αναμενόμενα ετήσια λειτουργικά οφέλη (\textlatin{OPEX savings}), λαμβάνοντας υπόψη τόσο τη μείωση της κατανάλωσης ηλεκτρικής ενέργειας όσο και την αποφυγή κόστους που συνδέεται με συμβάντα λειτουργικής αστοχίας και συντήρησης του εξοπλισμού ψύξης.

Το \textlatin{CAPEX} περιλαμβάνει το κόστος του υλικού εξοπλισμού, της εγκατάστασης, της παραμετροποίησης, της ολοκλήρωσης των επιμέρους υποσυστημάτων και της ανάπτυξης του λογισμικού που απαιτήθηκε για την επιχειρησιακή ωρίμανση της πλατφόρμας. Τα \textlatin{OPEX savings} εκτιμώνται με βάση τη μείωση της κατανάλωσης ενέργειας, χρησιμοποιώντας σταθερή τιμή αναφοράς 0{,}162~\euro/\textlatin{kWh}, βάσει των επίσημων λογαριασμών ηλεκτρικής ενέργειας της περιόδου αξιολόγησης, καθώς και με βάση τη μείωση των απωλειών που σχετίζονται με δυσλειτουργίες του εξοπλισμού και επαναλαμβανόμενα συμβάντα μη αποδοτικής λειτουργίας.

Σε αντίθεση με το τοπικό \textlatin{MVP}, το οποίο εξυπηρετούσε κυρίως τη συλλογή και την αρχική επιβεβαίωση της μεθοδολογίας, η \textlatin{Enterprise Cloud SaaS} πλατφόρμα παρείχε το κατάλληλο περιβάλλον για την παραγωγή προβολών \textlatin{ROI}, την απεικόνιση των ταμειακών ροών και τη διαμόρφωση ευανάγνωστων αναφορών προς τον τελικό χρήστη. Στο ίδιο πλαίσιο εντάχθηκε και η χρήση \textlatin{AI Energy Forecasting}, η οποία επέτρεψε την υποστήριξη προβλεπτικών σεναρίων και την ενίσχυση της λήψης αποφάσεων σε τεχνικό και διοικητικό επίπεδο.

Η τελική αξιολόγηση στηρίχθηκε σε δείκτες όπως ο χρόνος απόσβεσης (\textlatin{Payback Period}) και η πενταετής απόδοση επένδυσης (\textlatin{5-year ROI}), ώστε η αποτελεσματικότητα της πλατφόρμας να τεκμηριωθεί όχι μόνο σε τεχνικό αλλά και σε επιχειρησιακό επίπεδο. Με τον τρόπο αυτό, η μεθοδολογία του κεφαλαίου συνδέει άμεσα τη μετρητική διάσταση του ψηφιακού διδύμου με την πρακτική αξιοποίησή του ως εργαλείο βελτιστοποίησης και στρατηγικού σχεδιασμού. Το Σχήμα~\ref{fig:methodology_flow} παρουσιάζει το συνολικό διάγραμμα ροής της μεθοδολογίας αξιολόγησης.

\begin{figure}[htbp]
\centering
\includegraphics[width=\textwidth]{fig_5_1}
\caption{Διάγραμμα ροής της μεθοδολογίας αξιολόγησης και ενδεικτική απεικόνιση του περιβάλλοντος τεχνοοικονομικής ανάλυσης}
\label{fig:methodology_flow}
\end{figure}




